% =====================================================================
%  1 -- THE ASTROPHYSICAL SETTING.
%
%  Source: tier2.md Part 0 -- header (193-203), 0.1 (206-392),
%          0.2 (395-551), 0.3 (554-739), 0.5 (839-870),
%          0.5b (874-967), 0.6 (971-1009), 0.7 (1011-1023),
%          0.8 self-check (1025-1046).
%  Excluded and re-homed: 0.4 (743-837, "why conservation is
%          structural") opens section 2, immediately before the algebra
%          it is a statement about.  Nothing else moves.
%  Two ASCII diagrams (450-460, 613-618) are redrawn as TikZ figures.
% =====================================================================

\section{The astrophysical setting}
\label{part:setting}

Tier~0 Part~I asked \emph{why the project exists}. This section asks a narrower
and more technical question: \textbf{why does the algebra of Tier~2 take the
shape it does?} Why is there a cluster solver at all? Why is the central design
bet ``most reactions are in near-balance''? Why is \Ye{} the quantity every gate
in \code{CLAUDE.md} is written around, rather than, say, the \nuc{Ni}{56} mass
fraction?

None of that is derivable from the code. It comes from sixty years of silicon
burning theory, and the code is unreadable as \emph{physics} without it.

% =====================================================================
\subsection{The equilibrium hierarchy}
\label{sec:hierarchy}

\subsubsection{Three regimes, stated as claims about fluxes}
\label{sec:regimes}

Take the network ODE (Tier~0 \S II.4)
%
\begin{equation}
\dot Y_i \;=\; \sum_j \nu_{ij} R_j(Y; T, \rho),
\label{eq:networkode}
\end{equation}
%
and pair every strong\,/\,EM reaction $j$ with its inverse $\bar\jmath$
(Tier~1 \S III). Then the whole of nuclear astrophysics's equilibrium taxonomy
is a statement about which of the differences $\netflux_j = R_j - R_{\bar\jmath}$
are negligible.

\begin{center}
\begin{tabularx}{\textwidth}{@{}L{2.5cm} R{1.25} R{0.95} R{0.8}@{}}
\toprule
\textbf{regime} & \textbf{claim} & \textbf{composition parameters
  \emph{before} constraints} & \textbf{\ldots of which still free once
  $(T, \rho, \Ye)$ is given} \\
\midrule
\textbf{full kinetics} & nothing is assumed balanced &
  $n$ (all abundances) & $n - 2$ \\
\textbf{QSE} (quasi-statistical equilibrium) &
  reactions \emph{within} each of $g$ clusters balance; reactions
  \emph{between} clusters do not & $2 + g$ &
  \textbf{$g$} (the group masses) \\
\textbf{NSE} (nuclear statistical equilibrium) &
  every strong\,/\,EM reaction balances & 2 & \textbf{0} \\
\bottomrule
\end{tabularx}
\end{center}

Read the two right-hand columns together, because the familiar ``NSE has two
free parameters'' is shorthand for the third column and the \emph{point} of NSE
is the fourth: $(\mu_p, \mu_n)$ parameterise the composition, and the two
constraints ($\sum X = 1$, \Ye) then fix them, leaving nothing free. QSE's $g$
leftover degrees of freedom are exactly the group masses, which is why they have
to come from the slow dynamics.

The counting is the whole point, and it deserves to be derived rather than
asserted.

\paragraph{NSE $\to$ 2.} If every strong\,/\,EM reaction is balanced, then for
each such reaction the chemical potentials of reactants and products are equal
(\S\ref{sec:ensemble}--\S\ref{sec:binding}). Because every nucleus is reachable
from free nucleons by some chain of \emph{strong} reactions, the balance
conditions collapse to
%
\begin{equation}
\mu_i \;=\; Z_i\,\mu_p + N_i\,\mu_n \quad\text{for every } i.
\label{eq:nsemu}
\end{equation}
%
Two numbers, $(\mu_p, \mu_n)$, determine all $n$ abundances. Fixing them
requires two constraints: $\sum_i X_i = 1$ and $\sum_i (Z_i/A_i)X_i = \Ye$. So
NSE composition is a \emph{function}, $X(T, \rho, \Ye)$ --- no history, no
rates, no network. This is why NSE is so seductive computationally and why
\code{qse/solver.py} exists at all.

\paragraph{The reachability step is measurable, and the right measurement is not
the one you would reach for first.} ``Every nucleus is reachable from free
nucleons'' is a connectivity property of the \emph{strong\,/\,EM} subgraph, so
\code{graph\_metrics}' statement that the full exported graph is connected
\resultsref{2026-07-08} does not establish it --- the full graph carries weak
edges too. The evidence that does is \S\ref{sec:rank}'s: the left null space of
$\nu$ \textbf{restricted to strong\,/\,EM columns} is exactly 2-dimensional and
equals $\mathrm{span}\{A, Z\}$. A strong subgraph that fell into $c$
disconnected pieces would conserve each piece's mass separately and push that
nullity above 2; the measured 2 is precisely the statement this collapse needs.

\paragraph{QSE $\to$ $2 + g$.} Now suppose the network's reaction graph,
restricted to \emph{fast} reactions, is disconnected: it falls into $g$
clusters, plus the free nucleons. Within cluster $G$ the balance conditions
still collapse, but only relative to a cluster-local reference. The general
solution is
%
\begin{equation}
\mu_i \;=\; Z_i\,\mu_p + N_i\,\mu_n + \uG \quad\text{for } i \in G,
\label{eq:qsemu}
\end{equation}
%
with one extra offset \uG{} per cluster --- an integration constant of the
cluster's internal equilibrium that the \emph{slow} inter-cluster reactions must
determine. Setting all $\uG = 0$ recovers NSE. This is exactly the structure in
\code{qse/solver.py}'s docstring, and it is why the QSE solve has three unknowns
$(\up, \un, \uG)$ with a third constraint: the measured group mass fraction
$\sum_{i \in G} X_i$.

That third constraint is worth pausing on, because it is where QSE stops being a
closed theory. NSE needs $(T, \rho, \Ye)$ and nothing else. QSE needs
$(T, \rho, \Ye)$ \textbf{plus the group mass} --- a dynamical quantity that only
the slow bridge reactions can supply. QSE is therefore a \emph{reduction}, not a
solution: it turns an $n$-dimensional ODE into a $(g{+}1)$-dimensional one. That
is the entire content of the Hix \& Thielemann reduced-network programme, and it
is the structural ancestor of this project's Target-A design.

\paragraph{Full kinetics $\to$ $n$.} No reduction. What MESA\,/\,bbq actually
integrates, and what the emulator has to reproduce.

% ---------------------------------------------------------------------
\subsubsection{The historical arc --- and why each step mattered}
\label{sec:history}

This is worth knowing properly, because the project's design choices are
reactions to specific historical findings, not to a generic ``networks are
slow'' complaint.

\begin{itemize}
\item \textbf{\citet{hoyle1946}} (MNRAS 106, 343, ``The Synthesis of the
  Elements from Hydrogen'') --- the equilibrium calculation B$^2$FH later
  systematised and named the \emph{e-process}. The iron-peak abundance pattern
  is not the record of a reaction sequence; it is a thermodynamic equilibrium
  distribution, with \nuc{Fe}{56} the most abundant heavy nucleus in his
  freeze-out table (Table XIV). This is the origin of the idea that composition
  can be a \emph{state function}. The quantitative placement --- the
  \nuc{Fe}{56} peak at $\Ye = 26/56 \approx 0.464$ --- is not in Hoyle, whose
  neutron-richness variable is a degeneracy parameter; it comes from the later
  NSE surveys \citep[Fig.~2;][]{hwe1985}, the survey lineage beginning with
  \citet{clifford1965}.

  \begin{warn}
  \textbf{Say this one carefully.} The usual shorthand --- ``\nuc{Fe}{56}
  maximises binding energy per nucleon'' --- is not true: \nuc{Ni}{62} does
  (8.7946 vs 8.7904\,MeV/A, AME2020, \citealp{ame2020}). What NSE actually
  minimises is the free energy at fixed ($\sum X = 1$, \Ye), and near
  $\Ye = 0.464 = 26/56$ the winner is the nucleus whose $Z/A$ \emph{matches}
  \Ye{} while having near-maximal $B/A$ --- which is \nuc{Fe}{56} (the selection
  rule stated verbatim in \citealp{seitenzahl2008}). The \Ye{} qualifier is not
  decoration; it is the whole selection rule, and \S\ref{sec:fepeak} is the same
  statement read as a function of \Ye.
  \end{warn}

\item \textbf{\citet{bbfh1957}} (Rev. Mod. Phys. 29, 547) --- B$^2$FH
  systematises the e-process alongside s, r, p. Still equilibrium-first.

\item \textbf{\citet{bcf1968apjs}} (ApJS 16, 299, ``Nuclear Quasi-Equilibrium
  during Silicon Burning''; announced in \citealp{bcf1968prl}) --- the decisive
  correction, and the paper this tier is really downstream of. Full NSE is
  \emph{not} reached in silicon burning on the available timescale. Instead the
  network organises into \textbf{quasi-equilibrium}: in BCF's own picture a
  \emph{single} cluster, $28 \le A \le 62$, whose abundances are expressed
  relative to the slowly-depleting base nucleus \nuc{Si}{28} and the free light
  particles; the cluster is \emph{not} in equilibrium with the $A < 28$ sector,
  and the burning rate is carried by a few slow non-equilibrium links (chiefly
  \nuc{Mg}{24}$(\gamma,\alpha)$\nuc{Ne}{20}). Composition is neither a free ODE
  trajectory nor a state function; it is a state function \emph{of the group
  masses}, which evolve slowly.

\item \textbf{\citet{wac1973}} (ApJS 26, 231, ``The Explosive Burning of Oxygen
  and Silicon'') --- large explicit silicon-burning networks; QSE confirmed as
  an emergent property of the kinetics rather than an imposed assumption (``the
  striking tendency of elements to segregate into cluster-equilibrium groups'',
  their p.~250). This is also where the \textbf{two-group} refinement of BCF's
  single cluster enters: a silicon group ($A \le 45$) and an iron group
  ($A \ge 46$), joined by bridge reactions dominated by
  \nuc{Sc}{45}$(p,\gamma)$\nuc{Ti}{46} (their \S VIb, Fig.~17).

\item \textbf{\citet{ht1996}} (ApJ 460, 869; the reference cited in
  \code{configs/qse\_groups.yaml}) \textbf{and \citeyear{ht1999a}a}
  \citep[ApJ 511, 862;][]{ht1999a} --- the modern diagnostic study: a full
  299-species network run against QSE predictions, confirming the two groups
  split ``roughly A=24 to A=45'' (the boundary giving way to an $N = 23$ rule
  from Ti onward) and their merger into one cluster as burning proceeds. The
  $24 \le A < 45$ single-silicon-cluster span in the repo's default group config
  is a crisp paraphrase of this hedged boundary. The QSE \emph{reduced} networks
  themselves --- solve for $(\up, \un, \uG)$ plus the group masses, integrate
  only the bridges --- were made operational in the companion line of work:
  \citet{hix1998} and \citet{hix2007}.

\item \textbf{Guidry and collaborators, 2012--2013} \citep{guidry2012,
  guidry2013pe, guidry2013asym, guidry2013qss} --- \emph{partial equilibrium}
  and \emph{asymptotic} methods: rather than imposing group membership by
  $A$-range, detect equilibrated reaction pairs dynamically via a departure
  criterion, and stabilise explicit integration by algebraically removing them.
  The repo's $\varepsilon$-sweep criterion $|y - \bar y|/\bar y < \varepsilon$
  (\code{qse/diagnostics.delta\_species}) is this lineage --- it is Eq.~(8b) of
  the partial-equilibrium paper \citep{guidry2013pe}, whose own operative
  tolerance $\varepsilon = 0.01$ is the middle value of the repo's sweep.

\item \textbf{Machine-learning era} --- NNN \citep{grichener2025} and NuGNN
  \citep{kim2026}: learn the map directly. The physics-structural question these
  inherited but did not resolve is precisely BCF68\,/\,WAC73's: \emph{is the net
  flow carried by a small identifiable set of bridges?} If yes, a flux-space
  target with an equilibrium mask is the right architecture. If no, it is a
  liability.
\end{itemize}

\textbf{That question, restated as a measurement, is the kill-test}
(\S\ref{sec:killtest}). Tier~2 builds every instrument it needs.

% ---------------------------------------------------------------------
\subsubsection{What the project bets on this hierarchy}
\label{sec:bet}

Target A predicts a per-reaction flux \Phifl{} and maps it through the fixed
$\nu$. Three things make that attractive, and all three are QSE-dependent:

\begin{enumerate}
\item \textbf{The mask.} If most columns are equilibrated, the model only has to
  get a few dozen bridge fluxes right. That is a far easier learning problem
  than $n$ coupled abundances.
\item \textbf{Sparsity of the \emph{net} flow.} Even without a mask, if
  $|\netflux|$ is concentrated in $\lesssim 20$ columns, a loss weighted toward
  those columns is well-posed.
\item \textbf{Physical interpretability of errors.} A flux error is attributable
  to a named reaction; an abundance error is not.
\end{enumerate}

Every one of those is a claim about the \emph{data}, not about the architecture.
The honest posture --- and the one the repo takes --- is that Target A is a
hypothesis with a falsifier attached. \S\ref{sec:killtest} states the falsifier;
Tier~3 reports the verdict.

\begin{result}
\textbf{Read this before Tier 3:} the measured answer is not a clean yes or no.
The bridge flow \emph{is} concentrated (top-20 columns carry 0.82--0.88 of
inter-group flux, \resultsref{2026-07-12}) --- the BCF68\,/\,WAC73 picture
survives. But the \emph{maskable} set on the label manifold is \textbf{empty at
every $\varepsilon$} \resultsref{2026-07-12}, because the shipped labels' own
equilibria are displaced by a MESA bug (Tier~3, S13). Two different claims; the
data separates them. \S\ref{sec:negresults} carries the measurement.
\end{result}

% =====================================================================
\subsection{Silicon burning, mechanically}
\label{sec:sibmech}

\subsubsection{Why $\sim$3\,GK is the threshold}
\label{sec:threegk}

Tier~1 \S1.4 derived the photodisintegration rate: a reverse $(\gamma, \alpha)$
rate carries the factor $\exp(-Q/kT)$ relative to its forward
$(\alpha, \gamma)$ partner, times a phase-space ratio $\propto T^{3/2}$ (the
detailed-balance reciprocity relation: \citealp{fcz1967}; in modern form
\citealp{rauscher2000}, eq. for $\lambda_\gamma$). With $Q \approx 5$--8\,MeV
for the $\alpha$-captures in the Si--Ca range (AME2020: 6.95, 6.64, 7.04, 5.13,
7.70\,MeV for \nuc{Si}{28}$\to$\nuc{Cr}{48}) and $kT = 86.2\,\mathrm{keV}
\times \Tn$,
%
\begin{equation}
\frac{Q}{kT} \;=\; \frac{Q\,[\mathrm{MeV}]}{0.0862\,\Tn}.
\label{eq:qkt}
\end{equation}
%
At $\Tn = 2$ and $Q = 7$\,MeV that exponent is $\approx 41$ --- reverse rates
are suppressed by $e^{-41} \approx 10^{-18}$, utterly negligible; the flow is a
one-way ladder. At $\Tn = 4$ it is $\approx 20$, i.e.
$e^{-20} \approx 2\times10^{-9}$. At $\Tn = 6$ it is $\approx 13.5$.

\textbf{The exponent alone does not give the crossover, and it is worth doing
the comparison honestly}, because the shorthand ``$e^{-20}$ against a $10^{-9}$
abundance factor, so they become comparable'' quietly drops the largest term.
The quantity that decides is the ratio of \emph{fluxes}, not of rates:
%
\begin{equation}
\frac{\text{reverse flux}}{\text{forward flux}}
\;=\; \frac{\lambda_\gamma\,Y_B}
           {\rho N_A\langle\sigma v\rangle\,Y_A Y_\alpha},
\qquad
\frac{\lambda_\gamma}{\rho N_A\langle\sigma v\rangle}
\;=\; \frac{9.87\times10^{9}}{\rho}\,\Tn^{3/2}\,\Theta\,e^{-Q/kT},
\label{eq:fluxratio}
\end{equation}
%
with the constant $9.8685\times10^{9}$\,mol\,cm$^{-3}$ from
\citet{rauscher2000} and $\Theta$ the $O(1)$--$O(\text{few})$ ratio of
spin\,/\,partition-function\,/\,reduced-mass factors. At $\rho = 10^{8}$ and
$\Tn = 4$ the prefactor is ${\sim}10^{2}$--$10^{3}$, so the rate ratio is
${\sim}10^{-6}$ rather than the $2\times10^{-9}$ the exponent alone suggests
--- and \emph{that} is what the light-particle abundance factor
$Y_\alpha/Y_B \sim 10^{-6}$ meets. Two factors, both large, pulling the same
way; the qualitative conclusion below survives, but it does not follow from
$e^{-Q/kT}$ on its own.\derivedhere

The crossover is sharp because the exponent moves by ${\sim}Q/kT$ per unit
relative change in $T$: dropping from $\Tn = 4$ to 3 raises $Q/kT$ from 20 to
27, i.e. suppresses the reverse by $e^{-7} \approx 10^{-3}$. \textbf{The entire
regime change from ``ladder'' to ``equilibrium'' happens across less than a
factor of two in $T$} --- which is why the regime box ($\Tn$ 1.6--7.9) spans
both, why QSE onset is quoted at ${\sim}3$--3.3\,GK (\citealp{ht1999a} measure
quasi-equilibrium ``beginning to fail'' by $\Tn = 3.25$ with photodisintegration
freezeout near $\Tn \approx 3$, and NSE reached for $\Tn \gtrsim 5$; silicon
ignition itself is at 2.7--3.5\,GK, \citealp{woosleyjanka2005}), and why the
kill-test priority window is 3.3--5\,GK: that is where the structure is neither
absent (cold) nor total (NSE).

% ---------------------------------------------------------------------
\subsubsection{Two clusters, one set of bridges}
\label{sec:twoclusters}

Once photodisintegration is fast, the picture is Figure~\ref{fig:clusters}.

\begin{figure}[H]
\centering
\begin{tikzpicture}[
  font=\small,
  grp/.style={draw, line width=0.7pt, rounded corners=2pt,
              text width=5.4cm, inner sep=6pt, align=left,
              fill=resultrule!5},
  res/.style={draw, line width=0.7pt, rounded corners=2pt,
              inner sep=5pt, fill=exrule!6}]

\node[grp] (si) at (0,0)
  {\textbf{SILICON GROUP}\\[2pt]
   $24 \le A \le 44$\\[1pt]
   \nuc{Si}{28} \nuc{Si}{29} \nuc{Si}{30} \nuc{P}{31} \nuc{S}{32}
   \ldots\ \nuc{Ca}{40} \nuc{Ti}{44}\\[2pt]
   \emph{internally equilibrated}};

\node[grp] (fe) at (8.4,0)
  {\textbf{IRON GROUP}\\[2pt]
   $45 \le A \lesssim 60$\\[1pt]
   \nuc{Sc}{45} \ldots\ \nuc{Fe}{52} \nuc{Fe}{54} \nuc{Ni}{56} \ldots\\[2pt]
   \emph{internally equilibrated}};

\draw[-{Latex[length=2.4mm]}, line width=0.9pt]
  (si.east) -- node[above, font=\footnotesize]{bridge}
               node[below, font=\footnotesize]{slow} (fe.west);

\node[res] (pool) at (4.2,-2.6)
  {free $n$, $p$, $\alpha$ reservoir};

\draw[{Latex[length=2mm]}-{Latex[length=2mm]}, line width=0.7pt]
  (si.south) -- (pool.west);
\draw[{Latex[length=2mm]}-{Latex[length=2mm]}, line width=0.7pt]
  (fe.south) -- (pool.east);
\node[font=\footnotesize, align=center] at (4.2,-1.55)
  {fast $(\gamma,\alpha)\,(\alpha,\gamma)\,(p,\gamma)\,(\gamma,p)\,
    (n,\gamma)\,(\gamma,n)$};
\end{tikzpicture}
\caption{Silicon burning as two internally-equilibrated groups exchanging
material through a shared light-particle reservoir and a small set of slow
bridge reactions. Redrawn from the source's ASCII diagram; the edges are drawn
for the \code{default} group variant.}
\label{fig:clusters}
\end{figure}

The edges are drawn for the \textbf{\code{default}} group variant, and they are
drawn to the \emph{nucleon}: \code{configs/qse\_groups.yaml} gives
\code{a\_max: 45} and the loader applies \code{A < a\_max}
(\code{qse/diagnostics.py}), so \code{default} is $24 \le A \le 44$ and
\nuc{Sc}{45} sits on the \textbf{iron} side. The \code{a24\_46} variant moves
the line one nucleon up and puts \nuc{Sc}{45} inside. Getting this off by one is
not cosmetic --- it is exactly the difference between
\nuc{Sc}{45}$(p,\gamma)$\nuc{Ti}{46} being a bridge and being an internal
reaction, which is the whole content of \S\ref{sec:groupboundary}.

The mechanism BCF68 identified (with the two-group structure drawn above being
WAC73's refinement of their single cluster --- \S\ref{sec:history}):
\nuc{Si}{28} does not fuse with \nuc{Si}{28} (the Coulomb barrier at $Z = 14$ on
$Z = 14$ is prohibitive at these temperatures --- Tier~1 \S1.2.4). Instead,
\emph{some} \nuc{Si}{28} photodisintegrates, releasing $\alpha$, $p$, $n$ into a
reservoir; those light particles are then captured by \emph{other}
silicon-group nuclei, walking them upward. The net effect is
\nuc{Si}{28} $\to$ \nuc{Ni}{56}, but the mechanism is
disassembly-and-reassembly, not direct fusion. Hence ``rearrangement''.

Consequences that matter for the algebra downstream:

\begin{itemize}
\item \textbf{The light-particle reservoir is the coupling.} $n$, $p$, $\alpha$
  are the shared currency; they appear in nearly every column of $\nu$. This is
  why the isotope projection of the bipartite graph has radius 1 and diameter 2
  \resultsref{2026-07-08}: \emph{any} two species are two hops apart through the
  light particles. Message passing over this graph saturates almost immediately,
  which is both good ($K = 5$ suffices, \S\ref{sec:graphK}) and dangerous (a GNN
  can shortcut everything through $p$\,/\,$n$\,/\,$\alpha$ node features).
\item \textbf{The bridge set is small and identifiable.} BCF68's structural
  prediction, sharpened by WAC73 into the two-group form with a named dominant
  bridge. Measured, on the relaxed manifold: \msn{80}'s top-5 inter-group
  carriers are \code{Ne22(a,n)Mg25} (share 0.20), \code{Al27(p,a)Mg24} (0.13),
  \code{P31(p,a)Si28} (0.11), \code{Na23(a,p)Mg26}, \code{Mg24(n,g)Mg25};
  top-20 carry 0.88. \msn{151} adds the literature bottleneck
  \nuc{Sc}{45}$(p,\gamma)$\nuc{Ti}{46} \citep[\S VIb: ``always the dominant, or
  at least one of the dominant, flows linking the two groups''; reaffirmed by
  \citealp{ht1996}]{wac1973} at rank 2/251 under the boundary-at-46 group
  variant. \resultsref{2026-07-12}
\item \textbf{The group boundary is a modelling choice, not a fact.} Whether
  \nuc{Sc}{45} is ``in'' the silicon group changes which reactions count as
  bridges. The repo makes this a config variant
  (\code{configs/qse\_groups.yaml}: \code{default} $24 \le A < 45$ vs
  \code{a24\_46}) and requires both to be reported ---
  \S\ref{sec:groupboundary}.
\end{itemize}

% ---------------------------------------------------------------------
\subsubsection{The Fe-peak endpoint, and why \texorpdfstring{\Ye}{Ye} picks it}
\label{sec:fepeak}

At $\Ye = 0.5$ (equal protons and neutrons) the most-bound nucleus available is
\nuc{Ni}{56} ($Z = N = 28$, doubly magic --- Tier~0 \S0.3.2). Silicon burning at
$\Ye = 0.5$ terminates on \nuc{Ni}{56}, which later decays down the $A = 56$
chain (\nuc{Ni}{56} $\to$ \nuc{Co}{56} $\to$ \nuc{Fe}{56}, half-lives 6.075\,d
and 77.24\,d; \citealp{ensdf}) and powers the supernova light curve. Both steps
are weak, but not the same weak channel: \textbf{\nuc{Ni}{56} decays by electron
capture, essentially 100\%} --- it is not a $\beta^+$ emitter (the dominant
branch feeds the 1720\,keV level of \nuc{Co}{56}, leaving only ${\sim}0.4$\,MeV,
below the $2m_ec^2$ positron threshold; measured $\beta^+$ intensity
${\sim}10^{-3}$\,\%) --- while \nuc{Co}{56} is ${\sim}81\%$ EC and
${\sim}19\%$ $\beta^+$ \citep{ensdf}. Worth keeping straight given how carefully
\S\ref{sec:betaplus} separates the two ledgers.

But \Ye{} is not 0.5. Electron captures during silicon burning drive it down
(Tier~0 \S II.5), and at $\Ye < 0.5$ the equilibrium favours more neutron-rich
Fe-peak species: \nuc{Fe}{54}, \nuc{Ni}{58}, and at lower \Ye{} still,
\nuc{Fe}{56} itself --- each species peaking where its own $Z/A$ matches \Ye{}
\citep[Fig.~2: \nuc{Ni}{56} dominant for $\Ye \gtrsim 0.49$,
\nuc{Fe}{54}\,/\,\nuc{Ni}{58} for ${\approx}0.47$--0.49, \nuc{Fe}{56} below
${\approx}0.47$]{hwe1985}. The endpoint of silicon burning is therefore \emph{a
function of \Ye}, and the \nuc{Ni}{56} mass that eventually lights the supernova
is set by how much of the burnt material stayed near $\Ye = 0.5$.

\begin{invariant}
This is the physical reason the whole project is organised around \Ye{} rather
than around \nuc{Ni}{56}: \textbf{\Ye{} is the upstream variable, \nuc{Ni}{56}
the downstream consequence.} An emulator that gets \Ye{} right will get the
endpoint right for the right reason; one tuned on \nuc{Ni}{56} can be right
about \nuc{Ni}{56} and wrong about everything else.
\end{invariant}

% ---------------------------------------------------------------------
\subsubsection{Why the flow is not a ladder --- the trap for intuition}
\label{sec:notladder}

It is very tempting to picture silicon burning as \nuc{Si}{28} $\to$
\nuc{S}{32} $\to$ \nuc{Ar}{36} $\to$ \nuc{Ca}{40} $\to$ \nuc{Ti}{44} $\to$
\nuc{Cr}{48} $\to$ \nuc{Fe}{52} $\to$ \nuc{Ni}{56}, an $\alpha$-ladder with each
rung a net flux. Tier~0 \S0.6 already warns about this. Here is why it matters
\emph{algebraically}:

If the flow were a ladder, $\nu$ restricted to the flow-carrying columns would
be nearly bidiagonal, the net fluxes would be a chain of comparable magnitudes,
and \kap{} would be $O(1)$ on every carrying column. Target A would be trivially
well-posed.

What actually happens is that each ladder rung is a \emph{nearly balanced pair}
($\kap \ll 1$) whose small net difference is what advances the chain, while the
light-particle reservoir couples every rung to every other. The net flow through
the ladder is the \emph{residual} of large opposing gross fluxes --- which is
precisely the catastrophic-cancellation structure of Tier~0 \S V.1, and
precisely what makes \kap{} a condition number (\S\ref{sec:kappa-cond}).
\textbf{The ladder picture and the cancellation problem are the same fact seen
from two sides.}

% =====================================================================
\subsection{The \texorpdfstring{\Ye}{Ye} chain, derived end to end}
\label{sec:yechain}

This subsection is the ``theory behind the motivation''. Every gate in
\code{CLAUDE.md} is a \Ye{} gate; this is why.

\subsubsection{\texorpdfstring{$\Mch \propto \Ye^2$}{Mch proportional to Ye squared} --- the derivation}
\label{sec:mch}

Tier~0 \S I.2 does this; the short version, because \S\ref{sec:budget} needs the
constant.

A cold, fully degenerate, relativistic electron gas has pressure
%
\begin{equation}
P_e \;=\; K\,n_e^{4/3}, \qquad
K = \frac{\hbar c}{4}\left(3\pi^2\right)^{1/3},
\label{eq:degpressure}
\end{equation}
%
and $n_e = \rho \Ye / m_u$. So $P = K (\Ye/m_u)^{4/3} \rho^{4/3}$ --- a
polytrope of index $n = 3$. For an $n = 3$ polytrope the mass is
\emph{independent of central density}:
%
\begin{equation}
M \;=\; 4\pi\left(\frac{K'}{\pi G}\right)^{3/2}\!\!\cdot\,
        \bigl(-\xi_1^2\theta'(\xi_1)\bigr),
\qquad K' = K\,(\Ye/m_u)^{4/3},
\label{eq:polytropemass}
\end{equation}
%
with the Lane--Emden constant $-\xi_1^2\theta'(\xi_1) \approx 2.018$. Since
$M \propto K'^{3/2} \propto \Ye^2$,
%
\begin{equation}
\boxed{\;\Mch \;\approx\; 5.83\,\Ye^2\ \Msun.\;}
\label{eq:mch}
\end{equation}
%
This is the standard Chandrasekhar result \citep{chandrasekhar1931,
chandrasekhar1939}, usually written $\Mch = 1.457\,(2/\mu_e)^2\,\Msun$
\citep[eq.~3.3.17]{shapiro1983} --- identical with $\mu_e = 1/\Ye$. It is the
zero-temperature ideal-degenerate value; Coulomb corrections lower the actual
white-dwarf limit to $\approx 1.40\,\Msun$ \citep{hamada1961}, and general
relativity destabilises it slightly further \citep[\S6.10]{shapiro1983}, but the
$\Ye^2$ scaling --- all that the argument below uses --- survives.

At $\Ye = 0.5$, $\Mch = 1.46\,\Msun$; at $\Ye = 0.45$, $1.18\,\Msun$.
\textbf{A 10\% change in \Ye{} moves the Chandrasekhar mass by
$0.28\,\Msun$ --- about 19\%.}\derivedhere

% ---------------------------------------------------------------------
\subsubsection{Deleptonisation and the homologous core}
\label{sec:delepton}

The iron core grows by silicon shell burning until it exceeds \Mch{} and
collapses. During collapse, electron capture --- dominantly on Fe-peak
\emph{nuclei}, with free protons subdominant \citep{langanke2003, hix2003} ---
continues to lower \Ye: with modern capture rates the central \Ye{} falls to
$\approx 0.25$--0.27 after neutrino trapping (trapped lepton fraction
$Y_L \approx 0.285$--0.30), essentially frozen through bounce
\citep{janka2012}. The classic treatment \citep{bethe1990} quotes the trapped
\emph{lepton} fraction, $Y_L \approx 0.36$--0.39 --- note the two eras and the
\Ye-vs-$Y_L$ distinction; an earlier draft of this paragraph conflated them
(corrected in the 2026-08-14 audit). The \emph{inner} core --- the part that
collapses subsonically and homologously \citep{goldreich1980, yahil1983} --- has
a mass that scales with the same $\Mch \propto \Ye^2$ relation evaluated at the
trapped \Ye, and it is the inner core's mass that determines where the bounce
shock forms (``whose mass scales with the instantaneous Chandrasekhar mass
$M_{\rm Ch}(t) \propto \Ye^2(t)$, and whose edge defines the location of shock
formation at bounce'' --- \citealp{janka2012}).

Two chains, then, both keyed on \Ye{} (Figure~\ref{fig:yechain}).

\begin{figure}[H]
\centering
\begin{tikzpicture}[
  font=\small, node distance=3mm,
  bx/.style={draw, line width=0.6pt, rounded corners=2pt,
             inner sep=4pt, fill=resultrule!5},
  ar/.style={-{Latex[length=2mm]}, line width=0.6pt}]

\node[bx] (ye) {Si-burning \Ye};

\node[bx, right=14mm of ye, yshift=9mm] (core)
  {pre-collapse core mass \& structure};
\node[bx, right=6mm of core] (comp) {compactness};

\node[bx, right=14mm of ye, yshift=-9mm] (init) {initial \Ye{} for collapse};
\node[bx, right=6mm of init] (trap) {trapped \Ye};
\node[bx, right=6mm of trap] (minner) {$M_{\rm inner}$};

\node[bx, right=10mm of minner, yshift=7mm] (shock) {shock birth radius};
\node[bx, right=10mm of minner, yshift=-7mm] (expl) {explodability};

\draw[ar] (ye.east) -- (core.west);
\draw[ar] (core) -- (comp);
\draw[ar] (ye.east) -- (init.west);
\draw[ar] (init) -- (trap);
\draw[ar] (trap) -- (minner);
\draw[ar] (minner.east) -- (shock.west);
\draw[ar] (minner.east) -- (expl.west);
\end{tikzpicture}
\caption{The two chains keyed on \Ye. Redrawn from the source's ASCII chain.
The upper chain is why an \emph{emulator's} \Ye{} matters even before collapse.}
\label{fig:yechain}
\end{figure}

The first is why an \emph{emulator's} \Ye{} matters even before collapse: it
sets the initial condition of everything after.

% ---------------------------------------------------------------------
\subsubsection{The observable end}
\label{sec:observable}

The chain terminates in things telescopes see:

\begin{itemize}
\item \textbf{\nuc{Ni}{56} mass} --- sets the radioactive tail of the light
  curve, and in \emph{stripped-envelope} events (Ib/c) approximately the peak
  too, via Arnett's rule \citep[applied to SE-SN samples in e.g.
  \citealp{lyman2016}]{arnett1982}. \textbf{The Arnett-rule-at-peak statement
  does not transfer to the H-rich Type II events a 20\,\Msun{} progenitor
  typically produces}: there the plateau is powered by hydrogen recombination
  \citep{popov1993, kasen2009}, and \nuc{Ni}{56} sets the tail and modulates the
  plateau \emph{length}, not the peak luminosity \citep{kasen2009, nakar2016}.
  Either way, \Ye{} sets how much of the ejecta reaches \nuc{Ni}{56} rather than
  \nuc{Fe}{54}\,/\,\nuc{Ni}{58}.
\item \textbf{The Fe-peak isotopic ratios} --- \nuc{Ni}{57}/\nuc{Ni}{56},
  \nuc{Ni}{58}/\nuc{Fe}{56} --- measured directly in SN 1987A (\nuc{Co}{57}
  $\gamma$-lines: \citealp{kurfess1992}; \citealp{seitenzahl2014}) and in the
  solar abundance pattern, and directly sensitive to the neutron excess
  $\eta = 1 - 2\Ye$ \citep{wac1973, thielemann1996}.
\item \textbf{The explosion/failure branch} itself, hence the observed
  progenitor-mass--outcome mapping.
\end{itemize}

\foreignreg{0}{R-21} has the observational grounding in detail. The point for
Tier~2 is just that the chain is \emph{long and multiplicative}, which is what
makes the error budget in \S\ref{sec:budget} tight.

% ---------------------------------------------------------------------
\subsubsection{The error budget this implies}
\label{sec:budget}

Now the part that produces the numbers in \code{CLAUDE.md}.

Suppose the emulator makes a per-step error $\delta$ in \Ye, and a trajectory
has $N$ steps. Tier~0 \S IV.3 derives the accumulation trichotomy:

\begin{center}
\begin{tabular}{@{}ll@{}}
\toprule
\textbf{error character} & \textbf{accumulated after $N$ steps} \\
\midrule
systematic (biased)              & $N\cdot\delta$ \\
random walk                      & $\sqrt N\cdot\delta$ \\
contractive (stable attractor)   & $O(\delta)$ \\
\bottomrule
\end{tabular}
\end{center}

With $N \approx 1.6\times10^{3}$ steps and the \emph{conservative} systematic
assumption, demanding a total $|\dYe|$ below the physics floor of
${\sim}5\times10^{-3}$ gives
%
\begin{equation}
\delta \;\lesssim\; \frac{5\times10^{-3}}{1.6\times10^{3}}
        \;\approx\; 3\times10^{-6},
\label{eq:pergate}
\end{equation}
%
which is exactly the operative per-step gate in \code{CLAUDE.md}. The
$5\times10^{-3}$ floor itself is the FFN$\to$LMP weak-rate-library shift ---
\citet{heger2001} report that replacing the FFN rates \citep{ffn1980, ffn1985}
with the LMP shell-model rates \citep{lmp2000} raises the presupernova central
\Ye{} ``by $\dYe = 0.005$ to $0.015$'' --- i.e. \textbf{the gate is set by the
size of a known systematic uncertainty in the input physics.} That is the right
way to set it: an emulator whose error is smaller than the disagreement between
two defensible rate libraries is not the limiting factor.

Note what the gate is \emph{conditional on}: the assumption that error
accumulates linearly. \code{docs/phase0-checklist.md} requires the slope to be
measured; if it turns out to be $\sqrt N$, the gate can relax to
${\sim}5\times10^{-5}$. This is Tier~0 \S VII.3 --- decisions carry their own
reversal condition --- applied to the project's most load-bearing number. The
slope is unmeasured, which is why the gate is conditional: open, register
\reg{R-21}.

\begin{invariant}
\textbf{And this is why conservation is not negotiable.} A conservation
violation is \emph{by construction} systematic: it has a fixed sign per step
(\S\ref{sec:softpenalty}). It lands in the worst row of that table. Nothing else
in the error budget does. The theorem that makes conservation structural rather
than learned, and the arithmetic of the penalty alternative, open
\S\ref{part:conservation} at \S\ref{sec:whystructural}.
\end{invariant}

% ---------------------------------------------------------------------
\subsubsection{The output end --- and the derivative nobody has assembled}
\label{sec:outputend}

\S\ref{sec:mch}--\S\ref{sec:observable} walked $\Ye \to \Mch \to$ collapse $\to$
\nuc{Ni}{56} $\to$ light curve, and \S\ref{sec:budget} set the per-step gate.
Before leaving the chain it is worth being explicit about the end of it that is
\textbf{not} closed. \code{tier1.md} \S VIII.C.1 flags rate-uncertainty $\to$
\Ye{} propagation --- the \emph{input} end. This is the \emph{output} end, and
nothing flags it.

\S\ref{sec:yechain} derives $\Ye \to \Mch \to$ collapse $\to$ \nuc{Ni}{56} $\to$
light curve qualitatively, and \S\ref{sec:budget} sets the per-step gate from
the FFN$\to$LMP shift, $\dYe = 5\times10^{-3}$--$1.5\times10^{-2}$
\citep{heger2001}. Note what that is: \textbf{an input uncertainty used as an
output tolerance}. The logic --- ``an emulator whose error is smaller than the
disagreement between two defensible rate libraries is not the limiting factor''
--- is sound and is the right way to set a \emph{floor}. But it says nothing
about what error would actually change a scientific conclusion, because the
derivatives
%
\begin{equation}
\frac{\partial M(\nuc{Ni}{56})}{\partial \Ye},\qquad
\frac{\partial (\text{explodability})}{\partial \Ye},\qquad
\frac{\partial (\text{isotopic ratios})}{\partial \Ye}
\label{eq:outputderivs}
\end{equation}
%
have never been assembled --- not measured here, not cited from the literature,
not estimated. \S\ref{sec:mch} gets as far as $\dd\Mch/\dd\Ye$
($\Mch \approx 5.83\,\Ye^2\,\Msun$, so
$\dd\Mch/\dd\Ye = 2 \times 5.83\,\Ye = 11.7\,\Ye\,\Msun$, i.e.
\textbf{$5.8\,\Msun$ per unit $\delta\Ye$ at $\Ye = 0.5$}), and then the chain
stops.

Why this is worth closing:

\begin{itemize}
\item \textbf{It could relax or tighten the gate by orders of magnitude.} If
  explodability is insensitive to $\delta\Ye$ at $10^{-3}$, the $3\times10^{-6}$
  per-step gate is over-engineered by a factor of hundreds and the project is
  paying enormously for it (\S\ref{sec:costwall}'s cost wall is downstream of
  exactly this kind of demand). If some observable is sensitive at $10^{-4}$,
  the gate is right.
\item \textbf{It is the referee's first question.} ``You have built a
  $3\times10^{-6}$-per-step emulator; what observable changes at
  $3\times10^{-6}$?'' The current answer routes through an \emph{input}
  uncertainty, which is a defensible proxy but is not an answer.
\item \textbf{Much of it is sourceable rather than computable.} The CCSN
  literature has progenitor studies with \Ye{} variations; the derivative may be
  extractable from published model grids without running anything. That makes
  this a \emph{literature} task, which is cheaper than it sounds.
\end{itemize}

\code{tier1.md} \S VIII.E.7 (``residual sourcing gaps in the motivation chain'')
is the nearest existing item; this is the specific, actionable form of it.

Open (register: \reg{R-14}) --- it is the justification for the project's most
expensive requirement.

% =====================================================================
\subsection{What physically breaks when a leak happens}
\label{sec:leak}

Concretely, three things, in the order they bite.

\paragraph{1. The composition stops being a composition.}
$\sum X_i = \sum A_i Y_i = 1$ is not a normalisation convention you can
re-impose; it is the statement that the mass in the zone is the mass in the
zone. If $\sum_i A_i \dd Y_i = \varepsilon$ per step, then after $N$ steps
$\sum X_i = 1 + N\varepsilon$, and every downstream quantity that divides by it
--- mean molecular weight, $\Ye = \sum Z_i Y_i$, the mass density of any
species --- is wrong by that factor. Renormalising afterwards does not fix it:
renormalisation is a \emph{multiplicative} correction applied uniformly, whereas
the leak was species-specific. You have silently redistributed abundance between
species.

\paragraph{2. The EOS call downstream receives a state that does not exist.}
In a real stellar-evolution timestep the network's output composition is handed
to the equation of state, which computes $P(\rho, T, \text{composition})$ via
$\bar A$ and $\bar Z$ (Tier~0 \S0.1). A composition with $\sum X \neq 1$ makes
those means ill-defined. The hydro then integrates a pressure that corresponds
to no physical matter. This is the ``deployment coupling contract'' that Tier~1
\S VIII.E.6 flags as still unspecified --- but the \emph{requirement} is
unambiguous: whatever the coupling is, it will consume $\sum X = 1$.

\paragraph{3. \Ye{} drifts, so \Mch{} drifts, so the answer to the science
question drifts.} The charge row is the one with observable consequences.
\S\ref{sec:mch}: $\delta\Ye = 10^{-4}$ moves \Mch{} by
${\sim}6\times10^{-4}\,\Msun$; accumulated over a trajectory at $10^{-6}$/step
it is $1.6\times10^{-3}$ in \Ye{} and ${\sim}10^{-2}\,\Msun$ --- comparable to
the differences between progenitor models that people write papers about.

Which is the argument for the \emph{charge-to-lepton} form of the constraint
rather than plain charge conservation. \Ye{} \textbf{must} change --- that is
the physics. What must not happen is charge appearing or disappearing without a
lepton to account for it. \S\ref{sec:ledgers} builds the ledger that makes this
testable rather than tautological.

% =====================================================================
\subsection{Deployment --- what consumes the output, and what speedup is
            achievable}
\label{sec:deployment}

\S\ref{sec:leak} said what breaks when conservation fails, by naming three
downstream consumers. This subsection makes them explicit --- they are
\textbf{upstream of every accuracy target in the project} --- and then asks the
question that decides whether the project is worth doing at all: what speedup is
actually achievable. \code{tier1.md} \S VIII.E.6 flags the first half as
unspecified.

The unanswered question: \emph{what consumes the emulator's output, and what
does it need?}

Concretely, in a stellar-evolution timestep the network's output feeds:

\begin{center}
\begin{tabularx}{\textwidth}{@{}L{4.0cm} Y L{2.6cm}@{}}
\toprule
\textbf{consumer} & \textbf{needs} & \textbf{sensitivity} \\
\midrule
the EOS & $\bar A$, $\bar Z$ from the composition;
  \textbf{requires $\sum X = 1$} & \S\ref{sec:leak} \\
the energy equation & \enuc{} (and neutrino losses separately) &
  \S\ref{sec:energyidentity} \\
the next network call & $Y$ itself --- errors compound
  (\S\ref{sec:budget}) & Tier~0 \S IV.3 \\
the opacity\,/\,mixing modules & composition & unquantified \\
\bottomrule
\end{tabularx}
\end{center}

And the operator-split structure (Tier~0 \S IV.1) means the emulator's error
sits \emph{inside} an outer error the split already commits. Two things follow
that nobody has computed:

\begin{enumerate}
\item \textbf{Is the emulator's per-step error large or small compared with the
  operator split's own $O(\Delta t^2)$ commutator error?} If the split error
  dominates, a tighter emulator buys nothing, and the $3\times10^{-6}$ gate is
  over-specified. If it does not, the gate is the binding constraint. This is a
  one-experiment question --- integrate a zone with split and unsplit couplings
  and compare --- and it would either justify or relax the project's most
  load-bearing number.
\item \textbf{Is the coupling $(T, \rho)$-imposed or thermodynamically fed
  back?} \code{tier1.md} \S VIII.C.4 already flags that everything here is at
  \emph{imposed} $(T, \rho)$, whereas real burning heats the zone, which changes
  the rates. An emulator validated at imposed $(T, \rho)$ has not been validated
  in the loop it will run in, and the loop has a positive feedback (burning
  $\to$ heating $\to$ faster burning) that can amplify small errors.
\end{enumerate}

Both remain open (register: \reg{R-11}) --- the second is the reason the
accuracy gate has the value it has.

% ---------------------------------------------------------------------
\subsubsection{The cost ceiling --- why the fallback rate matters more than the
               speed}
\label{sec:costceiling}

Tier~0 \S I.3 and \S I.5 quantify the cost of the network and the speedup that
would matter; the Tier-4 sketch has \textbf{S22: UQ, OOD, and the
fallback-to-solver gate.} The arithmetic connecting them is one line, and it is
written nowhere.

Let $S_{\rm raw} = t_{\rm solver}/t_{\rm emu}$ be the raw speedup, $f$ the
fraction of calls that fall back to the real solver, and $t_{\rm det}$ the cost
of the out-of-distribution detector that decides. The emulator and the detector
are paid on \textbf{every} call; the solver is paid on the fallback fraction:
%
\begin{equation}
S \;=\; \frac{t_{\mathrm{solver}}}
             {t_{\mathrm{emu}} + t_{\mathrm{det}} + f\,t_{\mathrm{solver}}}
  \;=\; \Bigl(\tfrac{1}{S_{\mathrm{raw}}}
              + \tfrac{t_{\mathrm{det}}}{t_{\mathrm{solver}}} + f\Bigr)^{-1}
  \;\xrightarrow[\ S_{\mathrm{raw}}\to\infty\ ]{}\; \frac{1}{f}.
\label{eq:costceiling}
\end{equation}
\derivedhere

\textbf{The achievable speedup is capped at $1/f$ regardless of how fast the
emulator is.} Tabulated:

\begin{center}
\begin{tabular}{@{}ll@{}}
\toprule
\textbf{fallback rate $f$} & \textbf{ceiling on speedup} \\
\midrule
0.1\% & $1000\times$ \\
1\%   & $100\times$ \\
5\%   & $\mathbf{20\times}$ \\
20\%  & $\mathbf{5\times}$ \\
\bottomrule
\end{tabular}
\end{center}

Three consequences the project should be carrying and is not:

\begin{enumerate}
\item \textbf{S22's gate has no target.} A fallback-to-solver gate is specified
  with no stated acceptable fallback rate --- but the acceptable rate is
  \emph{derivable} from the speedup that would make the project worthwhile
  (Tier~0 \S I.5 has that number). $f_{\max} = 1/S_{\rm target}$. That single
  division turns an unspecified gate into a specified one.
\item \textbf{The detector is not free and is on the critical path.}
  $t_{\rm det}$ is paid on every call including the $1 - f$ that succeed. A
  detector costing 10\% of the solver caps the speedup at $10\times$ before any
  fallback happens. This makes ``expensive UQ'' (ensembles, MC-dropout at
  inference) structurally questionable, and cheap UQ (a single deterministic
  score) structurally preferred --- an architecture constraint arriving from
  economics rather than from statistics.
\item \textbf{Batching interacts badly with fallback.} In a stellar code the
  network is called per zone per timestep, and zones are naturally batched. If
  fallback is decided per zone, a batch containing one fallback zone pays both
  costs and may lose vectorisation on the rest. The effective $f$ is then closer
  to \emph{the probability that a batch contains any fallback} --- which for
  batch size $B$ and independent zones is $1 - (1-f)^B$, dramatically worse: at
  $f = 1\%$ and $B = 100$, 63\% of batches contain a fallback. Whether that
  matters depends entirely on the coupling contract (\S\ref{sec:deployment}),
  which is unspecified.
\end{enumerate}

Open (register: \reg{R-12}) --- arithmetic, an hour's work, and it converts two
unspecified design targets into numbers.

% =====================================================================
\subsection{The decision this tier serves}
\label{sec:killtest}

Everything in Tier~2 is instrumentation for one decision, and it is worth
writing the decision down as a linear-algebra question before building the
tools.

\begin{invariant}
\textbf{Kill-test.} Let $\mathcal{A} = \{j : \kap_j > 0.1\}$ be the \emph{active
set} --- the columns not in near-balance. Restricted to $\mathcal{A}$, is the
map $\netflux \mapsto \nu\netflux$ well-conditioned, and does it carry the
physics?
\end{invariant}

Operationally, from \code{CLAUDE.md}:

\begin{center}
\begin{tabularx}{\textwidth}{@{}Y L{1.4cm} L{1.6cm} L{4.6cm}@{}}
\toprule
\textbf{quantity} & \textbf{pass} & \textbf{fail} & \textbf{tool} \\
\midrule
coverage: fraction of dominant-isotope $|\Delta X|$ and of $|\dYe|$ carried by
  $\mathcal{A}$ & $\ge 95\%$ & --- &
  \S\ref{sec:kappa}, \code{killtest.active\_set.carried\_fractions} \\
spread: share of net columns near the \kap{} floor & --- & $> {\sim}30\%$ &
  \S\ref{sec:vacuous} \\
$\cond(\Sact) = \cond(\nu$ restricted to $\mathcal{A})$ & $< 10^{6}$ &
  $> 10^{8}$ & \S\ref{sec:rank}, \code{cond\_s\_active} \\
maskable-set size under the Guidry $\varepsilon$ sweep & non-empty &
  empty $\Rightarrow$ full-width & \S\ref{sec:departure} \\
\bottomrule
\end{tabularx}
\end{center}

If it passes, Target A with an equilibrium mask is the architecture. If it fails
--- either because the flow spreads over too many near-cancelled columns, or
because the restricted $\nu$ is ill-conditioned --- the project falls back to
Target B: predict $\dd Y$ directly in a linear space and project. Note in
advance that the fallback costs no expressive power: the two targets reach the
same set of increments (\S\ref{sec:reachable}, proved in
\S\ref{sec:equivalence}).

\textbf{Notice what each tool in this tier is for}, now that the question is
stated:

\begin{itemize}
\item \textbf{S7} gives $\nu$, $C$, $P$, and $\cond(\cdot)$ --- the
  \emph{conditioning} half of the question, and the fallback (Target B) if the
  answer is no. \S\ref{part:conservation}.
\item \textbf{S8} gives \kap{} and the coverage statistics --- the \emph{active
  set} half. \S\ref{part:fluxes}.
\item \textbf{S9} gives the independent NSE\,/\,QSE reference, without which
  ``in equilibrium'' has no measurable definition, and the Guidry $\delta$
  criterion that defines maskability. \S\ref{part:equilibrium}.
\item \textbf{S10} gives the reference integrator that produces the \Phifl{}
  labels Target A needs and that the shipped dataset does not contain.
  \S\ref{part:integration}.
\end{itemize}

Nothing in this tier is exploratory. Every module answers a clause of that one
question.

% =====================================================================
\subsection{What is \emph{not} in this tier}
\label{sec:scope}

To keep the boundaries clean:

\begin{itemize}
\item \textbf{Where $\lambda$ comes from} --- Tier~1. Tier~2 consumes $\lambda$
  and asks nothing about its provenance except that the pf gate held.
\item \textbf{Silicon burning as a physical process on real trajectories, the
  label pathology, the verdict} --- Tier~3 (S11--S16). This section gives the
  \emph{theory} of Si burning because the algebra needs it; the
  \emph{measurements} on actual burning histories are one tier later.
\item \textbf{The emulator} --- Tier~4, unwritten (\code{tier1.md} \S VIII.E.1).
\item \textbf{Screening and the two-\kap{} convention} --- Tier~1 \S IV. Tier~2
  uses the rule (equilibrium detection on unscreened \kap) without re-deriving
  the offset.
\end{itemize}

% =====================================================================
\subsection{Self-check}
\label{sec:selfcheck-setting}

Answer without scrolling up.

\begin{selfcheck}
\item NSE has 2 free parameters given $(T, \rho, \Ye)$. Where does the number 2
  come from, and what breaks the count for QSE?
\item Why does silicon burning proceed by disassembly-and-reassembly rather than
  by \nuc{Si}{28} $+$ \nuc{Si}{28}?
\item Write $\Mch \propto \Ye^2$ and say which step of the derivation makes the
  mass independent of central density.
\item State the theorem ``conservation laws are left null vectors'' including
  the quantifier that makes it load-bearing. \emph{(Answered in
  \S\ref{sec:leftnull}, which opens \S\ref{part:conservation}.)}
\item A soft conservation penalty achieves $\|C\,\dd Y\| \sim 10^{-8}$ per step.
  Work out what that accumulates to over $N = 1.6\times10^{3}$ steps and compare
  it against the \emph{right} budget --- then say why the penalty is still the
  wrong design in two independent ways that have nothing to do with that number.
  \emph{(Answered in \S\ref{sec:softpenalty}.)}
\item Why is the constraint \emph{charge-to-lepton closure} rather than charge
  conservation?
\item State the kill-test as a question about a matrix, and name which Tier-2
  node supplies each of its four measured quantities.
\end{selfcheck}
